\documentclass[10pt,letterpaper]{article}
\usepackage[margin=0.72in]{geometry}
\usepackage[T1]{fontenc}
\usepackage{lmodern}
\usepackage{microtype}
\usepackage{booktabs}
\usepackage{longtable}
\usepackage{array}
\usepackage{tabularx}
\usepackage{xcolor}
\usepackage{hyperref}
\usepackage{enumitem}
\usepackage{fancyvrb}
\usepackage{amsmath}
\usepackage{siunitx}
\hypersetup{colorlinks=true,linkcolor=blue,urlcolor=blue}
\setlist{nosep,leftmargin=1.5em}
\setlength{\parindent}{0pt}
\setlength{\parskip}{4pt}
\newcommand{\code}[1]{\texttt{\detokenize{#1}}}
\newcommand{\obs}{\textbf{Observed:} }
\newcommand{\infv}{\textbf{Inferred:} }
\newcommand{\hyp}{\textbf{Hypothesis:} }
\newcommand{\decision}{\textbf{Decision:} }

\title{\textbf{CAROM BridgeLearn Shadow-Mode v2}\\
Exact Experiment, Parameters, Complete Recorded Results, and Interpretation}
\author{ZeroBot research record for Benjamin Goertzel}
\date{22 July 2026}

\begin{document}
\maketitle

\begin{abstract}
This report documents the exact CAROM BridgeLearn v0.3.0 shadow-mode v2
calibration run.  The run observed all 1,500 AdamW updates at matched
per-update cadence while leaving the OneCycleLR schedule unchanged.  It
recorded parameter-variance transitions, split-batch noise estimates,
random-direction Hessian-vector curvature probes, sharpness forecasts, and
AR(2) dynamics for four parameter blocks.  The key positive result is that
one-step predicted and realized parameter variances agree almost perfectly
(correlations 0.999996--0.999999), showing that the large discrepancy in the
earlier cadence-mismatched run was mainly a horizon artifact.  The active
control gate nevertheless fails: predicted innovation is orders of magnitude
below empirical transition innovation, sharpness intervals are extremely
wide and near-zero-confidence, and the AdamW trajectory is stably second
order rather than adequately AR(1).  These failures are not all equivalent
to ordinary scalar hyperparameter mistuning.  Some parameters require
calibration, but the innovation mapping and state model require structural
changes before active BridgeLearn control is scientifically justified.
\end{abstract}

\tableofcontents

\section{Question and status}

The experiment asked:
\begin{quote}
Does matched per-update observation, Hessian-vector-product sharpness,
split-batch innovation measurement, and an AR(2) diagnostic make BridgeLearn
sufficiently calibrated for an active CAROM control experiment?
\end{quote}

\obs The command exited with status 0.  All 1,500 shadow non-mutation checks
passed, all telemetry is finite strict JSON, and the run produced 1,500
per-update observations, 30 split-batch measurements, 1,498 available AR(2)
fits, and 11 evaluation summaries.

\decision The mechanical shadow integration passes.  Activation does not:
BridgeLearn must remain in shadow mode until innovation, sharpness, and
momentum-state adequacy gates pass.

\section{Exact provenance and environment}

\begingroup\sloppy\small
\begin{tabularx}{\textwidth}{>{\raggedright\arraybackslash}p{0.27\textwidth}X}
\toprule
Item & Recorded value \\
\midrule
Experiment ID & \code{20260722T200000Z-bridgelearn-shadow-mode-v2} \\
Experiment directory &
\code{projects/carom/experiments/20260722T200000Z-bridgelearn-shadow-mode-v2/} \\
CAROM repository & \code{projects/carom/repos/carom} \\
Workspace branch & \code{agent/chat-room-identity-phase1} \\
Recorded workspace commit & \code{ffdc2934e8006a27ed9e0f7d0e7137f8456833fd} \\
Runner status & Runner and experiment were uncommitted in a shared dirty
workspace; exact runner content is preserved by SHA-256. \\
Runner SHA-256 &
\code{9443bb973b7885131aaf6c3506403c408e3b427d5e10e17b5b04f105e38ac2c5} \\
BridgeLearn & Version 0.3.0, dedicated local virtual environment \\
Python & 3.10.12 \\
PyTorch & 2.13.0+cpu \\
Platform & Linux x86\_64, CPU only \\
Elapsed wall time & 703.929 s in telemetry; 11:46.49 from \code{/usr/bin/time -v} \\
CPU utilization & 461 percent aggregate \\
Maximum resident set & 1,008,852 KiB \\
Exit status & 0 \\
\bottomrule
\end{tabularx}
\endgroup

This was not a paid or remote run.  No RunPod resource was active.

\section{Exact command}

\begin{Verbatim}[fontsize=\small]
cd /home/openclaw/research-agent/projects/carom/repos/carom
/usr/bin/time -v \
  /home/openclaw/research-agent/scratch/bridgelearn-carom-venv/bin/python \
  run_carom_bridge_shadow_v2.py scheduled \
  --steps 1500 \
  --seed 0 \
  --device cpu \
  --output \
  /home/openclaw/research-agent/projects/carom/experiments/\
20260722T200000Z-bridgelearn-shadow-mode-v2/telemetry.json
\end{Verbatim}

All unspecified command-line arguments took the runner defaults listed below.

\section{Complete parameter specification}

\subsection{Task and model}

\begin{longtable}{p{0.34\textwidth}p{0.58\textwidth}}
\toprule
Parameter & Exact value or behavior \\
\midrule
\endhead
Execution variant & \code{scheduled}; commands are applied sequentially in
presentation order. \\
Workspace & Six slots over $\mathbf{Z}_8$. \\
Program length & Uniform integer from 2 through 5 for each example; padded
with trailing no-ops to length 5. \\
Command vocabulary & Ten tokens: no-op, increment, double, negate, reverse,
shift-right, swap-first-two, cumulative modular sum, increment synonym, and
reverse synonym. \\
Input generation & Procedural; six independently uniform values in
$\{0,\ldots,7\}$ followed by exact execution of the sampled command list. \\
Model width & $d=48$. \\
Operator count & $K=8$. \\
Routing & Learned command embedding with $K$ logits, softmax temperature
$\tau=1.0$. \\
Shared operator core & Eight parallel read-transform-gate-write operators,
with attention, GELU transform, sigmoid gate, and residual workspace update. \\
Embeddings & Symbol embedding for nine values including blank; six position
embeddings. \\
Readout & LayerNorm, linear $48\rightarrow96$, GELU, linear
$96\rightarrow8$. \\
Parameter blocks & embeddings; routing; shared operator core; readout plus
any otherwise unassigned parameters. \\
Training loss & Cross-entropy over all six output slots. \\
Training batch size & 128. \\
Evaluation batch size & 512, newly generated from the same procedural
distribution. \\
Seed & 0 for Python \code{random}, NumPy, and PyTorch. \\
HVP random seed & 173 (training seed plus 173). \\
\bottomrule
\end{longtable}

\subsection{Optimizer and schedule}

\begin{longtable}{p{0.34\textwidth}p{0.58\textwidth}}
\toprule
Parameter & Exact value or behavior \\
\midrule
\endhead
Optimizer & PyTorch AdamW. \\
Learning-rate argument / OneCycle maximum & $2.0\times10^{-3}$. \\
OneCycleLR unspecified defaults & \code{pct_start=0.3},
\code{anneal_strategy='cos'}, \code{cycle_momentum=True},
\code{base_momentum=0.85}, \code{max_momentum=0.95},
\code{div_factor=25}, \code{final_div_factor=10000},
\code{three_phase=False}. \\
Initial LR & $8.0\times10^{-5}$. \\
Peak LR & Approximately $2.0\times10^{-3}$ near update 450. \\
Final recorded pre-update LR & $8.0\times10^{-9}$ at update 1499. \\
AdamW defaults not overridden & $\beta_1=0.9$, $\beta_2=0.999$,
$\epsilon=10^{-8}$, AMSGrad false, maximize false. \\
Weight decay & $10^{-4}$. \\
Gradient clipping & Global norm clipped to 1.0 before each optimizer update. \\
Updates & 1,500. \\
Evaluation cadence & Every 150 updates and at final update. \\
\bottomrule
\end{longtable}

\subsection{BridgeLearn observer and planner parameters}

\begin{longtable}{p{0.38\textwidth}p{0.54\textwidth}}
\toprule
Component & Exact parameters \\
\midrule
\endhead
Controller mode & \code{shadow}; proposed controls were recorded but not
applied to optimizer or scheduler state. \\
Observation cadence & Before and after every optimizer update. \\
Empirical transition & Paired pre/post-update flattened parameter samples,
separately for four blocks. \\
Split-batch noise & Batch divided into two halves every 50 updates; floor
$10^{-12}$; exponential filter $\alpha=0.2$; latest estimate reused between
measurements. \\
Sharpness probe & One independent random unit vector per block and update;
absolute Rayleigh quotient $|v^\top H v|$ computed with autograd HVP. \\
Sharpness filter & Exponential filter $\alpha=0.1$. \\
Sharpness plant & Default BridgeLearn \code{SharpnessPlant}; consecutive
filtered HVP values; progress signal fixed to one. \\
AR(2) batch observer & Rolling window 128 updates; 16 evenly spaced fixed
parameter coordinates per block; minimum 24 samples. \\
AR(2) online cross-check & Forgetting-factor recursive least squares with
forgetting 0.99. \\
AR(1) adequacy threshold & 0.5. \\
Initial curvature & 1.0 per block. \\
Initial noise scale & $10^{-8}$ per block. \\
Initial AR(1) adequacy & 0.5 per block. \\
Commanded contraction & $1-\eta h(\eta)$. \\
Commanded innovation &
$\max(\eta^2\,\mathrm{noise}/\mathrm{effective\ batch},10^{-15})$. \\
Reference filter & $\alpha=0.03$; maximum contraction change 0.01; maximum
log-innovation change 0.20. \\
Step actuator bounds & $10^{-7}$ to $5\times10^{-3}$; maximum log-step change
0.15. \\
Batch actuator & Fixed at 128; quantum 128. \\
Planner terminal variance & Half the initial block variance, floored at
$10^{-12}$. \\
Planner horizon & 1,500; maximum horizon multiplier 16; approximate progress
permitted. \\
Probe policy & Confidence threshold 0.4; cold variance 0.01; maximum probe
innovation $10^{-7}$; maximum total probe $10^{-5}$. \\
Momentum handling & First-order momentum permitted; AR(1) and structural
diagnostics recorded but not required for shadow planning. \\
\bottomrule
\end{longtable}

\section{Evaluation trajectory}

\begin{center}
\small
\begin{tabular}{rrrrr}
\toprule
Update & LR & Train loss & Eval accuracy & Max $|$variance error$|$ \\
\midrule
0    & 8.000e-5 & 2.1001 & 0.1175 & 2.270e-5 \\
150  & 5.619e-4 & 2.0440 & 0.1810 & 1.916e-4 \\
300  & 1.524e-3 & 1.9784 & 0.2012 & 4.148e-4 \\
450  & 2.000e-3 & 1.8199 & 0.2497 & 1.055e-3 \\
600  & 1.900e-3 & 1.7245 & 0.3171 & 8.795e-4 \\
750  & 1.621e-3 & 1.5875 & 0.3275 & 4.829e-4 \\
900  & 1.220e-3 & 1.5415 & 0.3730 & 3.937e-4 \\
1050 & 7.746e-4 & 1.4008 & 0.4141 & 2.586e-4 \\
1200 & 3.742e-4 & 1.4253 & 0.4346 & 1.613e-4 \\
1350 & 9.774e-5 & 1.3387 & 0.4424 & 3.259e-5 \\
1499 & 8.000e-9 & 1.4378 & 0.4137 & 1.475e-8 \\
\bottomrule
\end{tabular}
\end{center}

Across all 1,500 stochastic training batches, mean loss was 1.66668, minimum
loss was 1.19589, mean training accuracy was 0.31809, and maximum single-batch
training accuracy was 0.52474.  These summaries are descriptive only: the run
was designed to calibrate observers, not to estimate final model efficacy.

\section{Transition and variance calibration}

\begin{center}
\small
\begin{tabular}{lrrrrr}
\toprule
Block & Bias & MAE & RMSE & Max abs. error & Correlation \\
\midrule
Embeddings & $-4.63$e-5 & 6.58e-5 & 1.03e-4 & 5.33e-4 & 0.9999961 \\
Routing & $+3.74$e-4 & 3.74e-4 & 4.84e-4 & 1.24e-3 & 0.9999982 \\
Operator core & $+2.62$e-6 & 2.62e-6 & 3.62e-6 & 1.04e-5 & 0.9999988 \\
Readout & $+1.19$e-5 & 1.21e-5 & 1.60e-5 & 5.08e-5 & 0.9999988 \\
\bottomrule
\end{tabular}
\end{center}

The residual 5th, 50th, and 95th percentiles were:

\begin{center}
\small
\begin{tabular}{lrrr}
\toprule
Block & q05 & q50 & q95 \\
\midrule
Embeddings & $-2.27$e-4 & $-1.71$e-5 & $+6.05$e-5 \\
Routing & $+2.63$e-6 & $+3.03$e-4 & $+9.30$e-4 \\
Operator core & $+1.30$e-8 & $+1.96$e-6 & $+7.68$e-6 \\
Readout & $-6.01$e-7 & $+1.08$e-5 & $+3.13$e-5 \\
\bottomrule
\end{tabular}
\end{center}

Contraction MAEs for embeddings, routing, operator core, and readout were,
respectively, $3.19\times10^{-5}$, $1.76\times10^{-4}$,
$7.42\times10^{-4}$, and $2.68\times10^{-4}$.  The maximum variance error
occurred in routing at update 444, near the OneCycleLR peak.

\infv The earlier v1 maximum variance error of 0.111 and contraction gap of
0.338 were primarily caused by comparing one-step predictions with realized
statistics separated by 149--150 optimizer updates.  At matched horizon the
variance model is exceptionally accurate.  The small positive routing bias is
real in this trace but is not the main activation blocker.

\section{Innovation mismatch}

\begin{center}
\small
\begin{tabular}{lrrr}
\toprule
Block & Median predicted & Median empirical & Median split-batch unit noise \\
\midrule
Embeddings & 1.09e-12 & 9.25e-8 & 1.96e-4 \\
Routing & 2.07e-11 & 3.07e-7 & 6.17e-3 \\
Operator core & 5.11e-14 & 1.06e-7 & 1.15e-5 \\
Readout & 1.09e-11 & 6.62e-8 & 1.08e-3 \\
\bottomrule
\end{tabular}
\end{center}

Mean absolute predicted-versus-empirical innovation gaps were 2.25e-7,
7.52e-7, 2.16e-7, and 1.39e-7 in block order.

\obs The controller's current innovation estimate is several orders of
magnitude smaller than the empirical parameter-transition innovation.

\hyp Possible contributors are the current $\eta^2/B$ conversion, AdamW
preconditioning, momentum, weight decay, the split-batch estimator's scale,
and the fact that empirical innovation is inferred across parameter
coordinates rather than from replicated stochastic trajectories beginning at
the same parameter state.

\infv This is not yet identifiable as a single scalar tuning problem.  A
replicated same-state transition experiment is needed to determine whether a
constant scale correction suffices or whether the mapping must model optimizer
state and blockwise covariance.

\section{Sharpness estimates and forecast calibration}

\begin{center}
\small
\begin{tabular}{lrrrr}
\toprule
Block & Filtered HVP range & Center MAE & Median relative width & Final confidence \\
\midrule
Embeddings & 2.04e-5--6.99e-3 & 9.84e-5 & 73x & 5.47e-7 \\
Routing & 1.60e-6--9.77e-2 & 1.59e-3 & 3.70x & 5.47e-7 \\
Operator core & 4.94e-6--8.28e-4 & 1.81e-5 & 537x & 5.47e-7 \\
Readout & 1.88e-3--1.49e-2 & 2.27e-4 & 16.0x & 5.47e-7 \\
\bottomrule
\end{tabular}
\end{center}

Nominal empirical interval coverage was 99.8--100 percent.  This high coverage
is not evidence of a useful forecast because the intervals are extremely
wide and confidence is essentially zero.  The v1 gradient-secant values were
many orders larger than the v2 HVP values; the secant proxy should not be
treated as Hessian sharpness.

\hyp This component may admit parameter improvement---for example plant
warm-start, regularization, or better filter constants---but the single random
Rayleigh direction does not estimate the largest curvature eigenvalue.
Multiple HVP probes or a short power iteration is a structural estimator
upgrade, not merely parameter tuning.

\section{AR(2) dynamics and the AR(1) failure}

AR(1) adequacy was zero at threshold 0.5 for every block on every available
fit.  Mean adequacy scores were $1.18\times10^{-8}$,
$4.00\times10^{-10}$, $8.67\times10^{-9}$, and $6.41\times10^{-9}$.

\begin{center}
\small
\begin{tabular}{lrrrr}
\toprule
Block & Final $\phi_1$ & Final $\phi_2$ & Spectral radius & Stable-fit fraction \\
\midrule
Embeddings & 1.631 & -0.641 & 0.971 & 97.8\% \\
Routing & 1.818 & -0.823 & 0.971 & 90.5\% \\
Operator core & 1.567 & -0.579 & 0.970 & 100.0\% \\
Readout & 1.518 & -0.528 & 0.979 & 99.2\% \\
\bottomrule
\end{tabular}
\end{center}

The fitted dynamics are mostly stable under the Jury criterion; this is not a
finding of explosive optimization.  The negative second lag and spectral
radii just below one are consistent with persistent optimizer momentum.

\infv A scalar AR(1) model is structurally mismatched to this AdamW trajectory.
Tuning its coefficient cannot in general represent a stable AR(2)-like process.
A companion-state or explicit momentum adapter is the direct repair.

\section{What ``parameter tuning issue'' means here}

The data divide the problem into three categories:

\begin{enumerate}
\item \textbf{Already calibrated after measurement repair.}  One-step
parameter-variance prediction is highly accurate.  No major retuning is
indicated for this component, although routing has a small positive bias.

\item \textbf{Potential parameter and estimator calibration.}  Sharpness EMA,
plant warm-start, confidence regularization, number of HVP probes, and probe
policy thresholds should be swept in shadow mode.  Some are scalar parameters;
the move from one random direction to an eigenvalue-oriented estimate changes
the estimator itself.

\item \textbf{Structural mismatch before tuning.}  The innovation conversion
has not been calibrated against replicated same-state stochastic updates, and
the AR(1) state model omits AdamW momentum.  Sweeping controller gains on top
of these mismatches risks obtaining apparent calibration for the wrong reason.
\end{enumerate}

\decision The next experiment should remain shadow-only and should first
identify the innovation map and momentum state.  Only after those are fixed
should ordinary controller hyperparameters be tuned for active operation.

\section{Recommended bounded next gate}

\begin{enumerate}
\item Freeze selected parameter states at early, peak-LR, descending-LR, and
terminal schedule phases.
\item From each state, execute replicated stochastic updates with identical
parameters and optimizer state but independent batches.  Estimate blockwise
innovation covariance directly.
\item Fit and test candidate mappings: scalar rescale; $\eta^2/B$ with
AdamW-preconditioned gradient covariance; and a companion-state model that
includes first-moment momentum.
\item Compare AR(1), AR(2), and explicit AdamW-state predictions on held-out
update intervals using the same per-update cadence.
\item Replace the one-direction curvature probe with multiple probes or short
power iteration; tune EMA and plant regularization only after estimator
repeatability is quantified.
\item Require held-out calibration error, useful interval width, nontrivial
confidence, stability margins, and strict LR non-mutation before considering
active control.
\end{enumerate}

This sequence discriminates true parameter mistuning from model misspecification
without another expensive training run.

\section{Limitations}

\begin{itemize}
\item One seed and one CAROM scheduled variant were tested.
\item The training and evaluation data are procedurally sampled from the same
distribution; this is an observer-calibration run, not a generalization study.
\item The runner was uncommitted in a shared dirty workspace.  Its exact hash
is preserved, but commit-only reproduction is unavailable.
\item Sharpness uses one random direction per block per update and absolute
Rayleigh quotient; it is neither a trace estimate averaged over probes nor a
largest-eigenvalue estimate.
\item Split-batch and empirical innovations are not measurements of precisely
the same stochastic object.
\item AR(2) uses 16 fixed coordinates per block rather than the full parameter
state.
\item Evaluation points use fresh stochastic batches, so small endpoint
accuracy changes should not be overinterpreted.
\item BridgeLearn proposed controls were not applied.  This run provides no
evidence that active BridgeLearn control improves CAROM accuracy or stability.
\end{itemize}

\section{Telemetry field dictionary}

Every per-update record contains:
\begin{description}[style=nextline]
\item[\code{step}, \code{loss}, \code{training_accuracy}] Update index and
current stochastic-batch outcomes.
\item[\code{lr}, \code{lr_after_scheduler}] Four block learning rates before
and after the OneCycleLR scheduler step.
\item[\code{observed_variance}, \code{predicted_variance},
\code{calibration_error}] Per-block realized parameter population variance,
BridgeLearn forecast, and realized minus predicted value.
\item[\code{predicted_transition}, \code{realized_transition}] Per-block
contraction and innovation objects.
\item[\code{sharpness_raw}, \code{sharpness_estimate}] Raw and EMA-filtered
absolute random-direction HVP Rayleigh values.
\item[\code{sharpness_funnel}] Lower, center, upper, confidence, realized
error, and coverage for the forecast made on the prior update.
\item[\code{split_batch_innovation}] Per-block observation every 50 updates;
null on intervening updates.
\item[\code{ar2_adequacy}] Availability, $\phi_1$, $\phi_2$, innovation,
spectral radius, AR(1) residual autocorrelation, adequacy score, regressor
condition, Jury margins, stability verdict, and online RLS cross-check.
\item[\code{bridge_proposed_controls}] Controls proposed by BridgeLearn but
not applied.
\item[\code{shadow_lr_invariant}] Exact per-update Boolean that the shadow
apply did not mutate LR.
\item[\code{grad_norm}] Global gradient norm before clipping to 1.0.
\end{description}

The top-level telemetry also records schema and BridgeLearn versions, mode,
variant, seed, device, steps, batch sizes and cadences, block labels, elapsed
seconds, evaluation summaries, calibration report, control effort, final
evaluation accuracy, and serialized sharpness-plant state.  The complete
1,500-row, approximately 8.2 MB JSON is the authoritative raw-data appendix;
reprinting it verbatim inside the PDF would reduce usability without adding
information.

\section{Artifact verification}

\begin{longtable}{p{0.32\textwidth}p{0.62\textwidth}}
\toprule
Artifact & SHA-256 \\
\midrule
\endhead
\code{telemetry.json} &
\code{0b9b18d6b91452cdfc15928dea79ad7d39da3e2d3f23cdb7cc7460a32c25dc2c} \\
\code{stdout.log} &
\code{609c94ad70151409e54c17fc36debde6042d97b4ee19b04c0ccb14daab49af7c} \\
\code{stderr.log} &
\code{93cb210d703a0853e0ed14f7c4575c088a97dcff519184e35585631617bba6ec} \\
\code{command.sh} &
\code{618577204571c126f62e0a79d52ee8e03df9c0bfacc9b8a4f5b5cdc04fb28ead} \\
\code{run_carom_bridge_shadow_v2.py} &
\code{9443bb973b7885131aaf6c3506403c408e3b427d5e10e17b5b04f105e38ac2c5} \\
\bottomrule
\end{longtable}

\section{Bottom line}

\obs Matching the prediction and observation horizon removed the dramatic v1
variance anomaly.  BridgeLearn follows one-step block variance extremely well,
even near peak OneCycleLR.

\obs Innovation calibration, sharpness informativeness, and AR(1) adequacy
still fail.

\infv The remaining problem is partly tunable but not purely a conventional
hyperparameter issue.  The evidence specifically favors (i) direct
same-state innovation identification, (ii) an explicit momentum/AR(2)
state model, and (iii) a repeatable multi-probe curvature estimator, followed
by bounded shadow-mode tuning.

\decision Active control should not begin from the present scalar AR(1)
configuration.

\end{document}
